\section{Hardness of \tiling}
\label{sec:tiling}

We note that all of the problems
considered here are in $\NEXP$ since a valid tiling can be specified in
$2^{O(n)}$ bits and verified in time that is linear in the specification of the tiling.
For the result of the section, therefore, we focus on proving hardness.
The variants in Section~4 are also in $\NEXP$ by the same argument.

The construction will make use of a
binary counter Turing machine $M_{BC}$
which
starts with a blank semi-infinite tape. The head begins in a designated start
state in the left-most position of the tape. $M_{BC}$ will generate all
binary strings in lexicographic order. More specifically, there is a
function $f_{BC}: \mathbb{Z} \rightarrow \{0,1\}^*$ such that for some
constant $N_0$ and
every $N \ge N_0$, if $M_{BC}$ runs for $N$ steps, then the string $f_{BC}(N)$
will be written on the tape with the rest of the tape blank.
Moreover there are constants $c_1$ and $c_2$ such that
if $n$ is the length of the string $f_{BC}(N)$ and
$N \ge N_0$, then
$2^{c_1 n} \le N \le 2^{c_2 n}$.
We will also assume that for any binary string $x$, we can compute $N$
such that $f_{BC}(N) = x$ in time that is polynomial in the length of $x$.  In some of the variations of the problem we consider in
Section~4, we will need to put additional restrictions on $N$ (such as requiring $N$ to be odd), and in those cases, we still require
that we can find an $N$ with the appropriate restrictions such that $f_{BC}(N) = x$.

Using a standard padding argument, we can reduce any language in
$\NEXP$ to $\NTIME(2^{c_1 n})$. If $L$ is in $\NTIME(2^{n^k})$, the reduction
consists of padding an input $x$ so that its length is $|x|^k/c_1$~\cite{papa95}.
Thus, we will take an arbitrary non-deterministic Turing machine $M$ which accepts
an $\NEXP$-complete language $L$ in time $2^{c_1 n}$ and reduce it to \tiling.
The tiling rules and boundary conditions
will be specific to the Turing machine $M$ but will be independent
of any particular input. The reduction for \expref{Theorem}{th:unweighted-tiling}
then will take an input string $x$ and output
integer $N$ such that $f_{BC}(N-3)=x$. The tiling rules will have the property that
a string $x$ is in $L$ if and only if an $N \times N$ grid can be tiled according
to the tiling rules.
The Turing Machine $M_{BC}$ will be run for $N-3$ steps.
The $-3$ comes from the fact that the $N \times N$ grid holds two boundary rows as well
as the starting row.  This will leave the string $x$ on the final row of $M_{BC}$'s
computation.

\begin{proof}[Proof of \expref{Theorem}{th:unweighted-tiling}.]

The boundary conditions for the $N \times N$ grid will be that the four corners of the grid must have a designated tile type $\tileC$.
(We actually only need to use two corners as described in Section~4.1.)
First we will specify a set of boundary tiles and their constraints.
In addition to \tileC\,
there are four other kinds of boundary tiles: $\tileW$, $\tileN$, $\tileS$, $\tileE$.
We will call the rest of the tiles \emph{interior} tiles.
The tiling rules for the boundary tiles are summarized in \expref{Table}{table:fourcorners}.
\begin{table}
\begin{centering}
\begin{tabular}{lc|cccccc}
     &             &                \multicolumn{6}{c}{Tile on right} \\
     &             & $\tileC$ & $\tileW$ & $\tileE$ & $\tileN$ & $\tileS$ & $\tilevariable$ \\
 \hline
     & $\tileC$    &     N    &    N     &    N     &     Y    &     Y    &    N   \\
Tile & $\tileW$    &     N    &    N     &    N     &     N    &     N    &        \\
on   & $\tileE$    &     N    &    N     &    N     &     N    &     N    &    N   \\
left & $\tileN$    &     Y    &    N     &    N     &     Y    &     N    &    N   \\
     & $\tileS$    &     Y    &    N     &    N     &     N    &     Y    &    N   \\
 & $\tilevariable$ &     N    &    N     &          &     N    &     N    &        \\
\end{tabular}
\qquad
\begin{tabular}{lc|cccccc}
       &             &                \multicolumn{6}{c}{Tile on top} \\
       &             & $\tileC$ & $\tileW$ & $\tileE$ & $\tileN$ & $\tileS$ & $\tilevariable$ \\
 \hline
       & $\tileC$    &     N    &    Y     &    Y     &     N    &     N    &    N   \\
Tile   & $\tileW$    &     Y    &    Y     &    N     &     N    &     N    &    N   \\
on     & $\tileE$    &     Y    &    N     &    Y     &     N    &     N    &    N   \\
bottom & $\tileN$    &     N    &    N     &    N     &     N    &     N    &    N   \\
       & $\tileS$    &     N    &    N     &    N     &     N    &     N    &        \\
   & $\tilevariable$ &     N    &    N     &    N     &          &     N    &        \\
\end{tabular}
\caption{The tiling rules for boundary tiles.  $\tilevariable$ represents any interior tile.  ``N'' indicates a disallowed pairing of tiles.  For two boundary tiles, ``Y'' represents an allowed pairing.  Some of the rules for interior tiles are not specified, because they depend on the specific interior tile.}
\label{table:fourcorners}
\end{centering}
\end{table}

\tileW\ will mark the left side of the grid, \tileN\ the top of the
grid, \tileS\ the bottom of the grid, and \tileE\ the right side
of the grid. (See \expref{Figure}{fig:fourcorners}.)
\begin{figure}
\begin{centering}
\begin{tabular}{c@{\extracolsep{0.1em}}c@{}c@{}c@{}c}
\tileC & \tileN & \tileN & \tileN & \tileC \\
\tileW & & & & \tileE \\
\tileW & & & & \tileE \\
\tileW & & & & \tileE \\
\tileC & \tileS & \tileS & \tileS & \tileC
\end{tabular}
\caption{The only allowed tiling of the sides of a $5 \times 5$ grid.}
\label{fig:fourcorners}
\end{centering}
\end{figure}
To show that in any valid tiling the boundary tiles must be placed
in this manner, note the following facts:
\begin{itemize}
\item Nothing can go to the left of a $\tileW$ tile which means that
the only place a $\tileW$ tile could go is the left-most boundary.

\item Similarly, $\tileN$ tiles can only go in the
top row, $\tileS$ tiles can only go in the bottom row, and $\tileE$ tiles can only go in the right-most column.

\item No interior tile can border a $\tileC$ tile in any direction. Furthermore
a $\tileC$ cannot border on itself in any direction. This means that the
only possible locations for a $\tileC$ are the four corners since those
are the only places which can be surrounded by $\tileW$, $\tileN$, $\tileS$, or $\tileE$
tiles. Since the boundary conditions state that $\tileC$ tiles must go in
the corners, those are exactly the locations that will hold $\tileC$
tiles.

\item The only tiles that can go above or below a $\tileC$ tile are $\tileW$ and $\tileE$ tiles, and
$\tileE$ cannot go on the west boundary.  Thus, the tiles on the west boundary adjacent to the corners
must be $\tileW$ tiles.

\item Only $\tileW$ or $\tileC$ can go above or below a $\tileW$ tile, so the entire west boundary, except for the corners, will be $\tileW$ tiles.

\item Similar logic shows that the entire east boundary, except for the corners, will be $\tileE$ tiles.  Also, the entire south boundary, except for the corners, will be $\tileS$ tiles, and the entire north boundary, except for the corners, will be $\tileN$ tiles.

\end{itemize}

The remainder of the grid will be tiled in two layers.
The constraints on the two layers only interact at the bottom of the grid, so
we describe each layer separately. The actual type for an
interior tile is specified by a pair denoting its
layer 1 type and layer 2 type.
The bottom layer will be used to simulate the Turing machine $M_{BC}$.
The top boundary of the grid will be used to ensure that $M_{BC}$
begins with the proper initial conditions. Then the rules will enforce that
each row of the tiling going downwards advances the Turing machine $M_{BC}$
by one step. At the bottom of the grid, the output is copied onto layer 2.
Layer 2 is then used to simulate a generic non-deterministic Turing machine
on the input copied from layer 1. The lower left corner is used to
initialize the state of $M$ and the constraints enforce that
each row going upwards advances the Turing machine $M$ by one step.
Finally, the only states of $M$ that are allowed to be below an $\tileN$ tile
are accepting states.
Since each Turing machine only executes for $N-3$ steps and the grid has space
for $N-2$ tape symbols, the right end of the tape will never be reached.



Although it is well known that tiling rules are Turing complete~\cite{kmw62},
we review the ideas here in order to specify the details in our
construction. We will assume that the Turing machine $M$ is encoded
in a tiling that goes from bottom to top. This can easily be reversed for
$M_{BC}$ which goes from top to bottom.
The non-deterministic Turing machine $M$ is specified by
a triplet $(\Sigma, Q, \delta)$, with
designated blank symbol $\# \in \Sigma$, start state $q_0 \in Q$ and
accept state $q_A \in Q$.
There are three varieties of tiles, designated by elements of
$\Sigma$ (variety 1), $\Sigma \times Q \times \{r,l\}$ (variety 2) and
$\Sigma \times Q \times \{R,L\}$ (variety 3).  Variety 1 represents the state of the tape away from the Turing machine head.  Variety 2 represents the
state of the tape and head when the head has moved on to a location but before it has acted.  The $\{r, l\}$ symbol in a variety 2 tile tells us from which direction the head came in its last move.  Variety 3 represents the state of the tape and head after
the head has acted, and the $\{R,L\}$ symbol tells us which way the head moved.

The tiling rules are given in \expref{Table}{table:TM}.
%
\begin{table}
\begin{centering}
\begin{tabular}{lc|cccccc}
     &             &                \multicolumn{6}{c}{Tile on right} \\
     &             & $[b]$ & $[b,q',r]$ & $[b,q',l]$ & $[b,q',R]$ & $[b,q',L]$ & $\tileE$ \\
 \hline
     & $[a]$       &   Y   &     Y      &      N     &      Y     &     N      &   Y     \\
Tile & $[a,q,r]$   &   N   &     N      &      N     &      N     & If $q=q'$  &   N     \\
on   & $[a,q,l]$   &   Y   &     N      &      N     &      N     &     N      &   Y     \\
left & $[a,q,R]$   &   N   &     N      & If $q=q'$  &      N     &     N      &   N     \\
     & $[a,q,L]$   &   Y   &     N      &      N     &      N     &     N      &   Y     \\
     & $\tileW$    &   Y*  &     Y     & If $q'=q_0$ &      Y     &     N      &   Y     \\
\end{tabular}

\bigskip
\begin{tabular}{lc|ccccc}
       &             &                \multicolumn{5}{c}{Tile on top} \\
       &             & $[b]$    & $[b,q',r]$ & $[b,q',l]$                & $[b,q',R]$ & $[b,q',L]$ \\
 \hline
       & $[a]$       & If $a=b$ &  If $a=b$  &  If $a=b$, $q' \neq q_0$  &     N      &     N   \\
Tile   & $[a,q,r]$   &    N     &     N      &     N                     & If TM rule & If TM rule \\
on     & $[a,q,l]$   &    N     &     N      &     N                     & If TM rule & If TM rule \\
bottom & $[a,q,R]$   & If $a=b$ &  If $a=b$  &  If $a=b$, $q' \neq q_0$  &     N      &     N   \\
       & $[a,q,L]$   & If $a=b$ &  If $a=b$  &  If $a=b$, $q' \neq q_0$  &     N      &     N   \\
\end{tabular}
\caption{The tiling rules to simulate a Turing machine.  ``N'' indicates a disallowed pairing of tiles, and ``Y'' represents an allowed pairing.  Pairings with ``if'' statements are allowed only if the condition is satisfied.  ``If TM rule'' means the pairing is allowed only if $(a,q) \rightarrow (b,q',L/R)$ is one of the valid non-deterministic moves of the Turing machine being simulated.  The ``Y*'' entry will be modified later to get the correct starting configuration for the Turing machine. }
\label{table:TM}
\end{centering}
\end{table}
%
%
The table below gives an example of a section of tiles that
encodes the move $(a,q) \rightarrow (b,q',L)$ of the Turing machine from the bottom
row to the middle row and the move $(c,q') \rightarrow (f,q'',R)$ from
the middle row to the top row:

\[
\begin{array}{|c|c|c|}
\hline
[f,q'',R] & [b,q'',l] & [d] \\
\hline
[c,q',r] & [b,q',L] & [d] \\
\hline
[c] & [a,q,r] & [d,q,L]\\
\hline
\end{array}
\]

The lower row shows the head in the square with the $a$. The $[d,q,L]$
is from the previous TM move.
The tile $[b,q',L]$ in the middle row
enforces that the tiler is committing to executing the
step $(a,q) \rightarrow (b,q',L)$, although there may have been
other possible non-deterministic choices. The $[c,q',r]$ tile to the left of
the $[b,q',L]$ shows the new location and state of the head after the first move.
The $[b,q',L]$ tile now just acts as a $[b]$ tile for purposes of
the tiling above. The tile $[f,q'',R]$ enforces that the
tiler is committing to executing the step $(c,q') \rightarrow (f,q'',R)$.
The $[b,q'',l]$ tile to the right of
the $[f,q'',R]$ shows the new location and state of the head after the second move.


To be consistent with the horizontal tiling rules, a row must consist of some number of variety 1 tiles, along with adjacent pairs consisting of a variety 2 tile and a variety 3 tile.  The variety 2 and 3 tiles must be matched in the following sense: $q = q'$, and either we have $r$ on the left and $L$ on the right or $R$ on the left and $l$ on the right.
%
At the west edges of the row, we could possibly have just a single variety 2 tile by itself.
However, this can only happen if the tile is of the form $[a,q,l]$.
The rules enforce then that a \tileW\ tile can only go to the left of a variety 2 tile if
the Turing Machine state is $q_0$, which
is the starting head state of the machine.  We can assume without loss of
generality that the Turing machine never transitions back to the $q_0$ state, and
never has a transition that would move it left from the first location on the tape.
%
Given that we have one row of this form, the next row up must have the
same number of pairs of variety 2 and variety 3 tiles, and furthermore,
each pair must be shifted one position left or right, with the Turing
machine performing an allowed transition, as in the example.  If one row
has exactly one variety 2 tile $[a,q_0,l]$ in the leftmost position, then every row
above it in a valid tiling will also have exactly one variety 2 tile, and the $j$th
row up will correctly represent the state of the Turing machine tape and head after $j$ steps.

For our particular tiling, we will need additional rules for some pairs of boundary tiles and interior tiles.  These additional rules will be different for layers $1$ and $2$, and some of them will couple the two layers.  In addition, we will need to slightly modify some of the horizontal tiling rules.  The new and modified rules are summarized in \expref{Table}{table:TMboundaries}.  Also, for layer $1$, recall that we reverse ``above'' and ``below'' for the regular Turing machine simulation rules (\expref{Table}{table:TM}), so that time goes downwards instead of upwards.
\begin{table}
Additional layer $1$ rules:
\smallskip
\begin{centering}
\begin{tabular}{lc|ccccc}
\multicolumn{2}{l|}{Boundary Tile} &                \multicolumn{5}{c}{Adjacent interior tile} \\
\multicolumn{2}{l|}{Location} & $[a]$ & $[a,q,r]$ & $[a,q,l]$          & $[a,q,R]$ & $[a,q,L]$ \\
 \hline
Bottom & $\tileS$    &    Y        &     Y     &         Y             &    Y      &    Y     \\
Top    & $\tileN$    & If $a=\#$   &     N     & If $a=\#$ and $q=q_0$ &    N      &    N      \\
Left   & $\tileW$    & If $a \neq \#$ &  Y     & If $q=q_0$            &    Y      &    N      \\
\end{tabular}

\bigskip

\end{centering}
\noindent Additional layer $2$ rules:
\smallskip
\begin{centering}
\begin{tabular}{lc|ccccc}

\multicolumn{2}{l|}{Boundary Tile} &                \multicolumn{5}{c}{Adjacent interior tile} \\
\multicolumn{2}{l|}{Location} & $[a]$         & $[a,q,r]$  & $[a,q,l]$                            & $[a,q,R]$ & $[a,q,L]$ \\
 \hline
Bottom  & $\tileS$ & If $a$ matches layer $1$ &     N
& \parbox{30mm}{\centering If $q=q_0$ and\\ $a$ matches layer $1$} &    N      &    N      \\
Top     & $\tileN$ &   Y                      & If $q=q_A$ & If $q=q_A$                           &    Y      &    Y      \\
Left    & $\tileW$ & If $a \not\in \Sigma_{M_{BC}}$ &  Y   & If $q=q_0$                           &    Y      &    N      \\
\end{tabular}

\caption{Additional tiling rules between certain boundary and interior tiles for
layers $1$ and $2$.  Note that these rules are specific to the boundary
only because we have established that the boundary tiles can only
be placed at the boundary.
For layer $2$, the alphabet symbol $a \in \Sigma$ must match the alphabet symbol for the corresponding layer $1$ tile when on the bottom interior row.}
\label{table:TMboundaries}
\end{centering}
\end{table}

We would like to start out the
Turing machine $M_{BC}$ with $[\#, q_0, l]$ in the leftmost location
followed by $[\#]$ tiles.  This is ensured by the additional rules for layer $1$: The top interior row must consist of only these two types of tiles, and the leftmost location cannot be $[\#]$.  Since there cannot be any variety 3 tiles, there can only be one $[\#, q_0, l]$ in the leftmost location.
We can assume without loss of generality that
the Turing machine overwrites the leftmost $\#$ on the tape
and never writes a $\#$ there again.

The rest of the layer 1 rules just enforce the rules for the Turing machine
$M_{BC}$.

Now in order to copy the output from $M_{BC}$ to the input tape
for $M$, we restrict the kinds of tiles that can go above $\tileS$ tiles.
All the alphabet characters in the bottom interior row of layer 2 must
match the alphabet characters for layer 1.  That is, the output of $M_{BC}$ is
copied onto the input of $M$.

We also want to ensure that the starting configuration of $M$ has
only one head in the leftmost location.  To accomplish this,
we forbid a $\tileW$ to go next to an $[a]$ tile for $a \in \Sigma_{M_{BC}}$, so
the leftmost tile in the 2nd row from the bottom of layer 2 must be $[a,q_0,l]$.  Since there
can be no variety 3 tiles in the 2nd row, that must be the only variety 2 tile
in the row.  A little care must be taken to overwrite the leftmost input tape
character with something that is not in the alphabet of $M_{BC}$.
This is because we have forbidden having an $[a]$ tile to the right
of a $\tileW$ for any $a \in \Sigma_{M_{BC}}$, and this prohibition
applies to \emph{all} rows.
The information encoded in the left-most tape symbol can be retained by
having a new $a'$ symbol in $\Sigma_M$ for every $a \in \Sigma_{M_{BC}}$.

Finally, the only variety 2 tiles on layer 2 which we allow below a
$\tileN$ tile must be of the form $[a,q_A, r/l]$, where $q_A$ is the accepting state.
Thus, there is a valid tiling if and only if the non-deterministic TM $M$ accepts on input $x$ in $N-3$ steps.
\end{proof}
